%% ============================================================
%% PAPER 2: TOPOLOGY-, GEOMETRY- AND HEAT-CURVE-AWARE SIZING OF
%%          THERMAL STORAGE AND HEAT PUMPS IN DISTRICT HEATING
%% Target Journal: Energy Conversion and Management
%% Version: 0.1 — SCAFFOLD (placeholders only)
%% Framework: CALION
%% ------------------------------------------------------------
%% ECM constraints (Guide for Authors, retrieved 2026-07-20):
%%   - Full length article: max 9,000 words  -> target ~7,800
%%   - Max 15 figures + tables COMBINED in the main text
%%   - Highlights: 3-5 bullets, max 85 characters each (incl. spaces)
%%   - Graphical abstract: 531 x 1328 px (h x w), readable at 5 x 13 cm
%%   - Numbered sections (1, 1.1, 1.1.1); cross-refer by number
%%   - Numeric citations in square brackets
%%   - Data statement: Option C -> data MUST be deposited in a
%%     repository and cited (plan a Zenodo DOI for CALION)
%%   - Separate "Declaration of generative AI ..." section directly
%%     before the reference list
%%   - CRediT authorship contribution statement required
%% ============================================================
\documentclass[review,3p]{elsarticle}

\usepackage[utf8]{inputenc}
\usepackage[T1]{fontenc}
\usepackage{amsmath,amssymb,amsfonts}
\usepackage{graphicx}
\usepackage{booktabs}
\usepackage{multirow}
\usepackage{array}
\usepackage{siunitx}
\usepackage{xcolor}
\usepackage{hyperref}
\usepackage{cleveref}
\usepackage{lineno}
\usepackage{float}
\usepackage{subfig}
\usepackage{enumitem}

\modulolinenumbers[5]
\sisetup{per-mode=symbol}
\DeclareSIUnit{\eur}{EUR}
\DeclareSIUnit{\bar}{bar}
\hypersetup{colorlinks=true, linkcolor=blue, citecolor=blue, urlcolor=blue}
\biboptions{sort&compress}

%% --- Editorial helpers (remove before submission) ---
\newcommand{\todo}[1]{\textcolor{red}{[TODO: #1]}}
\newcommand{\placeholder}[1]{\textcolor{red}{[#1]}}
\newcommand{\wordbudget}[1]{\marginpar{\tiny\textcolor{gray}{#1 w}}}

%% --- Notation shared with Paper 1 (keep consistent!) ---
\newcommand{\Lthree}{\mathrm{L3}}
\newcommand{\Lplus}{\mathrm{L3}^{+}}
\newcommand{\kHC}{k}                      % heat curve slope parameter
\newcommand{\COP}{\mathrm{COP}}
\newcommand{\VTES}{V_{\mathrm{TES}}}
\newcommand{\QWP}{\dot{Q}_{\mathrm{HP}}}
\newcommand{\TAC}{\mathrm{TAC}}
\newcommand{\LCOH}{\mathrm{LCOH}}

\journal{Energy Conversion and Management}

\begin{document}

\begin{frontmatter}

%% ------------------------------------------------------------
%% TITLE / AUTHORS — handled by the author, left as placeholder
%% ------------------------------------------------------------
\title{\placeholder{Title -- to be set by author}}

\author[inst1]{\placeholder{Author list -- to be set by author}}

\affiliation[inst1]{organization={\placeholder{Affiliation}},
            country={Germany}}

%% ------------------------------------------------------------
%% ABSTRACT — write LAST (150-250 words)
%% ------------------------------------------------------------
\begin{abstract}
\todo{Write last, after Results and Discussion are final.
Structure: (i) problem -- heat-curve parameter, HP COP and TES capacity
are treated as separate fixed inputs in existing DH planning literature;
(ii) approach -- MILP with L3+ spatial hydraulics and McCormick
linearisation that optimises all three simultaneously, plus endogenous
siting; (iii) cases -- two real networks, deliberately contrasting;
(iv) headline numbers; (v) the topology-dependent siting asymmetry as
the transferable insight.}
\end{abstract}

%% ------------------------------------------------------------
%% HIGHLIGHTS — 3-5 bullets, MAX 85 CHARACTERS EACH. Write last.
%% ------------------------------------------------------------
\begin{highlights}
\item \todo{Bullet 1 -- simultaneous k/COP/V\_TES optimisation (<=85 chars)}
\item \todo{Bullet 2 -- Stadtbach headline result (<=85 chars)}
\item \todo{Bullet 3 -- Memmingen headline result (<=85 chars)}
\item \todo{Bullet 4 -- topology-dependent siting asymmetry (<=85 chars)}
\end{highlights}

%% ------------------------------------------------------------
\begin{keyword}
district heating \sep thermal energy storage \sep heat pump \sep
MILP \sep heat curve \sep endogenous siting
\end{keyword}

\end{frontmatter}

\linenumbers

%% ============================================================
\section{Introduction}
\label{sec:intro}
%% BUDGET: ~1,300 words | ELEMENTS: none
%% STATUS: writable now (no campaign data required)
%% SOURCES: Paper2_V1.bib; Paper 1 (MainEingereicht06072026.tex)
%% ------------------------------------------------------------
%% Intended arc:
%%   1.1 Motivation: electrification of industrial/municipal DH
%%   1.2 Related work: (a) TES sizing, (b) HP integration & siting,
%%       (c) supply-temperature / heat-curve optimisation,
%%       (d) network-resolved design MILPs -- and the gap: all three
%%       decision variables are coupled but never co-optimised
%%   1.3 Contribution & delimitation from Paper 1
%%   1.4 Paper outline
%% ------------------------------------------------------------
\todo{Section 1 -- drafted 2026-07-21. Citations verified against
Paper2\_V1.bib; two flagged below need author confirmation before
submission.}

\subsection{Motivation}
\label{sec:intro:motivation}

District heating is one of the few heat supply technologies that can
absorb electrification at the scale required by climate policy without
requiring parallel retrofits of every connected building, because the
transition happens once, centrally, at the generation side
\cite{Lund.2014,Werner.2017}. The transition is commonly framed as a
shift toward fourth- and fifth-generation systems built around lower
distribution temperatures, waste-heat recovery and demand-side
integration \cite{Lund.2018,Lund.2021,Buffa.2019,Gjoka.2023}, and it
increasingly means replacing fossil generation with large-scale heat
pumps, a shift already under way in several European district heating
markets \cite{Averfalk.2017}. Two
design questions come with it that fixed-fuel generation never had to
answer. First, a heat pump's efficiency and, in networks with volatile
sources such as industrial waste heat, its economic output depend on the
temperature lift it operates against -- so the network's own supply
temperature becomes a design lever rather than a fixed boundary
condition. Second, unlike a gas boiler, a heat pump cannot instantly
follow demand at low marginal cost, which makes thermal storage's
capacity and location a first-order economic decision rather than an
operational afterthought \cite{Christensen.2024}.

These two levers, together with the heat pump's own capacity and
location, are not independent. Lowering the network supply temperature
improves heat pump efficiency but, at fixed thermal capacity, inflates
the storage volume needed to hold a given quantity of heat -- the two
consequences of the same design decision pull the economics in opposite
directions. Where the storage and the heat pump are best placed depends
in turn on the network's topology and on which economic signal
dominates its operation, so the same network transformation stage can
imply different siting decisions in different systems. Existing district
heating planning practice generally treats supply temperature, heat
pump sizing and storage sizing as sequential, separately fixed inputs
\cite{Terlouw.2025}; Section~\ref{sec:intro:relatedwork} substantiates
this and locates the gap that motivates the present work.

\subsection{Related work}
\label{sec:intro:relatedwork}

The literature on district heating optimisation is extensive
\cite{Li.2017,Sameti.2017}; the six strands below are the ones that
bear directly on the coupling this paper addresses.

\emph{Thermal storage sizing.} A substantial literature addresses how
large a district heating storage should be and where its economic value
comes from, typically for a single, given network configuration
\cite{Dahash.2019,Dahash.2021,Sifnaios.2023,Renaldi.2017,Gabrielli.2018,
Ryan.2022,Lesko.2018,Nuytten.2013}. Within this strand, storage sited to
exploit demand-response and price flexibility has been studied
explicitly \cite{Capone.2021}, and storage capacity has been examined
specifically as a lever for absorbing industrial waste heat
\cite{Yang.2020} -- relevant background given Stadtbach's waste-heat
feed (Section~\ref{sec:cases:networks}). Storage sizing coupled to
network design has also been studied with an explicit topology: Ceruti
et al.\ \cite{Ceruti.2025} co-optimise network topology, generation
sizing and thermal storage
across synthetic and real networks under temporally resolved prices,
finding that large heat pumps and storage reinforce each other
economically -- the closest existing precedent to the present work.
Their formulation, however, holds the network's supply temperature
regime fixed and does not treat storage siting as a decision distinct
from network topology.

\emph{Heat pump integration and siting.} Where within a network a heat
pump should be installed, and how many, has been addressed directly:
network-resolved mixed-integer formulations optimise heat pump number,
capacity and location jointly with network operation
\cite{Hering.2021,Hering.2022,Capone.2023}, booster and central
configurations are compared explicitly \cite{Ostergaard.2016}, and
reviews catalogue the placement options and configurations available
for coupling heat pumps to district heating and cooling networks
\cite{BarcoBurgos.2022,Schlosser.2020,Jesper.2021}. A related strand
examines heat pump operation jointly with combined heat and power
\cite{Levihn.2017} and with flexible, price-responsive dispatch
\cite{Rahlf.2026,Abokersh.2020}. Separately, the accuracy consequences
of simplifying heat pump performance models -- the same methodological
question this paper faces in choosing a COP surrogate
(Section~\ref{sec:model:hp}) -- have been examined directly
\cite{Zaubitzer.2026,Kuosa.2024}. These studies establish
siting as a genuine network-resolved decision, as the present work also
treats it, but do not couple it to a storage siting decision or to a
variable supply-temperature regime.

\emph{Supply temperature and heat curve optimisation.} The economic
value of lowering distribution temperature has been examined at the
level of an individual heat-pump installation, where shifting the
heating curve within the outdoor-temperature range where it is
uneconomical to do so yields measurable savings \cite{Krug.2025}, and at
the level of a whole municipality's district heating share, where supply
temperature enters as a design variable alongside the heat source mix
\cite{Trentmann.2026}. Reviews and simulation studies quantify the
savings available from supply- and return-temperature reduction more
broadly \cite{Guelpa.2023,Wang.2024,Buonomano.2024}, including cascade
arrangements that connect a low-temperature branch to the return line of
a higher-temperature network \cite{Volkova.2020}. Li et al.\
\cite{Li.2022} come closest to the coupling exploited in the present
work, showing that the economically optimal distribution temperature
for a system with short-term storage differs from the lowest
technically feasible one -- but do not treat storage capacity or heat
pump siting as spatial, network-resolved decisions.

\emph{Demand-side flexibility and price-responsive operation.} A
separate literature studies how district heating can provide flexibility
to the electricity system rather than the other way around, through
demand response \cite{Guelpa.2021,Ju.2023}, thermal inertia
\cite{Guelpa.2022}, network-side balancing services
\cite{Boldrini.2022}, and integration with electricity markets more
broadly \cite{Golmohamadi.2022,Javanshir.2023}, with clustering and
operational strategies proposed to exploit this flexibility in practice
\cite{Guelpa.2018,Guelpa.2020}. Both networks studied here are exposed to
this kind of price signal in their operating cost
(Section~\ref{sec:cases:parameters}), but this strand addresses
operational flexibility rather than the investment-stage siting and
capacity decisions this paper is concerned with.

\emph{Network-resolved design optimisation.} A separate strand of work
formulates the district heating design problem -- pipe routing, node
selection, generation placement -- as a mixed-integer or mixed-integer
nonlinear program over the network graph
\cite{Mertz.2016,Bordin.2016,Blommaert.2020,Resimont.2021,Marty.2018,
Marquant.2018,Weinand.2019,Wang.2023,Wack.2023,Wack.2023b,Wack.2024,
Vieth.2025,Egberts.2020,Roland.2020,Guo.2024,Dorotic.2019,Sameti.2019}.
These formulations resolve topology and asset placement jointly and in
detail -- several address multi-period or multi-objective variants of
the same design problem -- but treat the operating temperature regime as
an exogenous input rather than a design variable, and generally do not
include thermal storage sizing as part of the same optimisation.

\emph{Adjacent operational and simulation frameworks.} A further,
largely separate literature develops detailed operational or simulation
frameworks for 4th- and 5th-generation networks -- MILP-based supply
temperature control \cite{Wirtz.2021}, decomposition approaches for
network planning \cite{Wirtz.2023}, dynamic pipe models
\cite{vanderHeijde.2017}, integrated design-and-control formulations
\cite{vanderHeijde.2019}, and general simulation environments
\cite{Schweiger.2017}. Closest in spirit to the present work's concern
with model fidelity, Wirtz et al.\ \cite{Wirtz.2021b} ask directly which
level of model complexity is sufficient for multi-energy system design
-- the same question Paper~1 \cite{Ruess.2026} answers for dispatch and
that Section~\ref{sec:model:framework} takes as its starting point. None
of this strand addresses the investment-and-siting coupling studied
here.

\emph{The gap.} Across these six strands, no formulation we are aware
of optimises heat pump siting and capacity, storage siting and capacity,
and the network's supply-temperature regime jointly, in a
network-resolved model, on real network topologies. Existing work
co-optimises at most two of these three axes at a time: topology with
storage \cite{Ceruti.2025}, siting with capacity
\cite{Hering.2021,Hering.2022}, or temperature with storage
\cite{Li.2022}. The present work closes this by treating all three
jointly, and does so on two real networks chosen specifically to differ
in geometry and scale, so that any result common to both is
attributable to the coupling itself rather than to a single network's
idiosyncrasies.

\subsection{Contribution and scope}
\label{sec:intro:contribution}

This paper builds directly on Paper~1 \cite{Ruess.2026}, which
establishes the dispatch-fidelity level $\Lplus$ used throughout this
work (Section~\ref{sec:model:framework}) but does not address investment
or siting. The present work takes $\Lplus$ as given and adds an
investment and siting layer on top of it: heat pump, electrode boiler
and thermal storage capacity and location become decision variables,
alongside a network transformation stage that governs the
supply-temperature regime.
\placeholder{Paper 1 is currently under review at Applied Energy; cite
as ``submitted'' and insert volume/page/DOI once accepted.}

The central methodological claim is not that these quantities are
optimised simultaneously inside a single mixed-integer program --
Section~\ref{sec:model:heatcurve} explains why the transformation stage
is treated as an enumerated set rather than a continuous decision
variable. Rather, the contribution is a nested
enumeration--optimization scheme: the network transformation stage is
enumerated exhaustively, and within each stage, heat pump, electrode
boiler and storage capacity and siting are optimised endogenously to
global optimality under a consistently propagated, temperature-dependent
coefficient of performance. Because every stage is solved to proven
optimality and the stage dimension is enumerated exhaustively rather
than sampled, the reported result is a global co-optimum over the full
product space of stages and investment decisions -- a stronger guarantee
than a single-shot mixed-integer program over a continuous heat-curve
parameter could offer, since, as Section~\ref{sec:model:heatcurve}
argues, such a parameter is not well defined independently of the
discrete retrofit programme each stage represents.

This paper addresses three research questions:
\begin{enumerate}[label=\textbf{RQ\arabic*.}, leftmargin=2.2em]
\item How does the network transformation stage interact with the
optimal siting and capacity of heat pump, electrode boiler and thermal
storage in a network-resolved investment model
(Sections~\ref{sec:results:sweep}--\ref{sec:results:coupling})?
\item Does a network's geometry and scale determine whether free siting
of heat pump and storage improves on the best available fixed-node
placement, or does a well-chosen fixed site already capture most of the
achievable benefit (Section~\ref{sec:results:siting})?
\item Once the storage-capacity penalty of a lower distribution
temperature -- reduced usable energy density, hence more volume per
unit of stored heat -- is priced in alongside the heat pump efficiency
gain it also confers, does reducing distribution temperature still
reduce total system cost, and by how much
(Sections~\ref{sec:results:coupling}--\ref{sec:discussion:topology})?
\end{enumerate}

\subsection{Outline}
\label{sec:intro:outline}

Section~\ref{sec:model} specifies the model: the balances and
constraints inherited from Paper~1, the network transformation stages,
the heat pump, electrode boiler and storage formulations, the investment
objective and its two linearisation techniques, and the validation
strategy. Section~\ref{sec:cases} introduces the two case-study networks
and the scenario matrix built on top of them.
Section~\ref{sec:results} reports the capacity--cost sweep, the
heat-curve--COP--storage coupling, and the fixed- and endogenous-siting
comparisons. Section~\ref{sec:discussion} examines the sensitivity of
the design to economic assumptions, explains why network topology
determines the siting strategy, and states the model's validation
results and limitations. Section~\ref{sec:conclusions} closes the paper.

%% ============================================================
\section{Model formulation}
\label{sec:model}
%% BUDGET: ~2,200 words | ELEMENTS: F1
%% STATUS: WRITABLE NOW -- fully number-independent
%% SOURCE: CALION_Paper2_Implementation_Statement.md (Parts A-C)
%% ------------------------------------------------------------
This section specifies the optimisation model. Section~\ref{sec:model:framework}
delimits what is inherited from the dispatch formulation of Paper~1
\cite{Ruess.2026} and what is added here; Sections~\ref{sec:model:balances}
to~\ref{sec:model:tes} present the physical and technological
constraints; Section~\ref{sec:model:objective} states the investment
layer, the objective function and the two linearisation techniques
employed; and Section~\ref{sec:model:validation} sets out the validation
strategy. \Cref{fig:F1} summarises the model architecture.

\begin{figure}[htbp]
  \centering
  \fbox{\parbox{0.85\linewidth}{\centering\vspace{3cm}
  TBD: model architecture diagram --- highest design priority per the
  project's figure specification; not yet produced.\vspace{3cm}}}
  \caption{Model architecture: network layer, investment layer and the
  nested enumeration--optimization scheme.}
  \label{fig:F1}
\end{figure}

\subsection{Modelling framework and fidelity level}
\label{sec:model:framework}

The model is implemented in the CALION framework and operates at the
fidelity level denoted $\Lplus$ in Paper~1 \cite{Ruess.2026}: a
node-and-pipe representation of the physical network in which supply
temperatures propagate spatially, thermal losses depend on the local
temperature level, pressure drop is represented by a linearised
Darcy--Weisbach relation, and transport delay is resolved at hourly
granularity. Paper~1 established this level by systematically comparing
five variants, from a copperplate abstraction without network physics to
a mixed-integer quadratically constrained formulation without
linearisation, and quantified both the topology-abstraction error and
the linearisation error incurred at $\Lplus$. That comparison is not
repeated here.

The present work takes $\Lplus$ as given and adds an investment and
siting layer on top of it. Concretely, the following elements are
inherited unchanged and are therefore stated only to the extent needed
for the model to be readable in isolation: the hydraulic representation,
the temperature-dependent loss propagation along pipes, the transport
delay formulation, and the general procedure by which a network is
calibrated against metered operational data. Of the two case-study
networks (Section~\ref{sec:cases}), Stadtbach is new to the
present work, while Memmingen is the same network used
as Paper~1's primary case (Section~\ref{sec:cases:networks}); its
calibration is nonetheless carried out here rather than simply
inherited, since Paper~2 additionally re-derived its topology and
re-cleaned its demand data.
\placeholder{Before submission: apply the anonymisation convention used
in Paper~1 (relative rather than absolute demand/cost figures) to this
network throughout -- Table~T1, T3/T4, and Figures~F3--F6 all currently
carry absolute numbers for it.}
Because Paper~2 additionally re-cleaned this network's demand and
topology data (Section~\ref{sec:cases:networks}), its calibration is
carried out here rather than simply inherited, even though the
underlying site is not new. The
following elements are new independently of the case studies: the
discrete capacity selection for thermal storage, the siting variables
and their colocation variants, the electrode boiler and heat pump
investment decisions, the network-transformation stages described in
Section~\ref{sec:model:heatcurve}, and the nested
enumeration--optimization scheme that links them.

A note on terminology is warranted. In what follows, an
\emph{investment variable} is a quantity the solver chooses within a
single optimisation run, whereas a \emph{stage} is a member of an
enumerated set over which independent runs are performed. Heat pump,
electrode boiler and storage capacity as well as siting are investment
variables. The network transformation stage is not.

This positioning -- an enumerated dimension traded against solver
tractability, rather than folding every design axis into a single
monolithic program -- follows general guidance on managing complexity in
energy system optimisation \cite{Kotzur.2021} and on choosing temporal
and structural aggregation levels that preserve design accuracy
\cite{Hoffmann.2022}; broader surveys of the energy-system modelling
landscape provide context for where the present tool sits within it
\cite{Hoffner.2024}.

\subsection{Nodal energy and mass balances}
\label{sec:model:balances}

For every node $i$ and hour $t$, thermal power is conserved:
\begin{equation}
\label{eq:heatbalance}
\sum_{a \in \mathcal{A}_i} q_{a,i}(t) + \sum_{j \in \mathcal{N}_i^{-}} \dot{Q}_{ji}(t)
= D_i(t) + \sum_{k \in \mathcal{N}_i^{+}} \dot{Q}_{ik}(t)
+ q_i^{\mathrm{ch}}(t) - q_i^{\mathrm{dis}}(t) + Q_i^{\mathrm{dump}}(t),
\end{equation}
where $\mathcal{A}_i$ is the set of generation assets at node $i$,
$\mathcal{N}_i^{-}$ and $\mathcal{N}_i^{+}$ are the sets of nodes
feeding into and fed from node~$i$, $D_i$ is heat demand,
$q_i^{\mathrm{ch}}$ and $q_i^{\mathrm{dis}}$ are storage charging and
discharging power, and $Q_i^{\mathrm{dump}}$ is a curtailment slack that
absorbs surplus waste heat. Network losses are embedded in the delivered
power $\dot{Q}_{jk}$ of each pipe rather than appearing as a separate
term.

Independently of \cref{eq:heatbalance}, water mass is conserved:
\begin{equation}
\label{eq:massbalance}
\sum_{j \in \mathcal{N}_i^{-}} \dot{m}_{ji}(t) + \dot{m}_i^{\mathrm{gen}}(t)
= \sum_{k \in \mathcal{N}_i^{+}} \dot{m}_{ik}(t) + \dot{m}_i^{\mathrm{dem}}(t),
\qquad
\dot{m}_i^{\mathrm{gen}}(t) = \frac{10^{3}\, q_i(t)}{c_p \left(T_{\mathrm{s},i} - T_{\mathrm{r},i}\right)} .
\end{equation}
Maintaining \cref{eq:massbalance} as a constraint in its own right,
rather than treating mass flow as a post-processed image of delivered
heat, matters at nodes that simultaneously host local generation and
pass flow downstream. At such nodes both the generation term and the
downstream delivery terms enter \cref{eq:heatbalance}, and their
interaction with \cref{eq:massbalance} is the point at which a
formulation that conserves energy globally can nevertheless violate mass
continuity locally. Both networks studied here contain such nodes.

\subsection{Network transformation stages and the heat curve}
\label{sec:model:heatcurve}

Supply temperature is governed by an ambient-compensated heat curve of
the standard parameterised form
\begin{equation}
\label{eq:heatcurve}
T_{\mathrm{s}}(t) = \min\!\left\{ T_{\mathrm{s}}^{\max},\,
\max\!\left\{ T_{\mathrm{s}}^{\min},\;
T_{\mathrm{s}}^{\min} + \kHC \left(T_{\mathrm{s}}^{\max} - T_{\mathrm{s}}^{\min}\right)
\frac{T_{\mathrm{h}} - T_{\mathrm{a}}(t)}{T_{\mathrm{h}} - T_{\mathrm{a}}^{\mathrm{des}}}
\right\} \right\},
\end{equation}
with heating limit temperature $T_{\mathrm{h}} = \SI{15}{\celsius}$,
design ambient temperature $T_{\mathrm{a}}^{\mathrm{des}} =
\SI{-12}{\celsius}$ -- both following DIN~EN~12831-1
\cite{DINEN12831.2017} and identical for the
two networks -- and slope parameter $\kHC$. Heat curve control of this type is standard practice and its
optimisation has been studied for individual heat pump installations
\cite{Krug.2025}; what is specific here is that the curve is applied at
network level and interacts with siting and capacity decisions.

Three stages, denoted HK0, HK1 and HK2, are considered. It is important
to be precise about what distinguishes them, because they are not three
values of a single parameter. Each stage bundles a slope $\kHC$, a
minimum supply temperature $T_{\mathrm{s}}^{\min}$ and, for the retrofit
stages, a reduced return temperature $T_{\mathrm{r}}$
(\cref{tab:T2}). HK0 ($\kHC = 1.0$; $T_{\mathrm{s}}^{\min} =
\SI{70}{\celsius}$ and \SI{74}{\celsius} for the two networks)
represents the existing building stock. HK1 ($\kHC = 0.8$;
\SI{65}{\celsius} and \SI{70}{\celsius}) represents retrofitting half of
the substations, which lowers the achievable return temperature. HK2
($\kHC = 0.6$; \SI{60}{\celsius} and \SI{66}{\celsius}) represents
complete substation retrofit. A fourth setting, denoted TVLFIX, holds
the supply temperature at its historical constant value and serves as
the reference against which heat curve control is assessed.

The stages are therefore alternative network-transformation narratives
rather than samples of a continuously adjustable design lever. This
distinction is not cosmetic: a continuous $\kHC$ would leave
$T_{\mathrm{s}}^{\min}$ and $T_{\mathrm{r}}$ undefined, since neither
follows from the slope, and both are consequences of a discrete
investment programme in the building stock. Treating the transformation
stage as an enumerated dimension is thus the modelling choice that
matches the underlying decision structure. It has the secondary benefit
of keeping $T_{\mathrm{s}}(t)$ exogenous within each run, which
preserves linearity in the constraints of
Sections~\ref{sec:model:hp} and~\ref{sec:model:tes}.

\subsection{Heat pump and electrode boiler models}
\label{sec:model:hp}

Heat pump electricity consumption is
\begin{equation}
\label{eq:pel}
P^{\mathrm{el}}(t) = \frac{\dot{Q}^{\mathrm{HR}}(t)}{\COP(t)}
 + \frac{\dot{Q}^{\mathrm{def}}(t)}{\COP^{\mathrm{def}}},
\end{equation}
where $\dot{Q}^{\mathrm{HR}}$ is heat recovered from the industrial
waste heat source and $\dot{Q}^{\mathrm{def}}$ is the share supplied at
a default performance figure when the source is unavailable. The
coefficient of performance follows a Lorenz-type correlation evaluated
on logarithmic mean temperatures on both the sink and the source side,
\begin{equation}
\label{eq:cop}
\COP(t) = A(t)\, B(t)\, \eta \left(1 - q_{\mathrm{ww}}(t)\right)
+ 1 - \eta - F_Q ,
\end{equation}
where $A(t)$ is the reciprocal Lorenz efficiency formed from the two
logarithmic mean temperatures, $B(t)$ corrects for pinch and internal
temperature losses, and $q_{\mathrm{ww}}$ and $F_Q$ are empirical loss
terms. The correlation assumes a source-side spread of \SI{20}{\kelvin}
and a pinch of \SI{5}{\kelvin} on both heat exchangers, and the result is
clipped to $[1.01, 12]$. It is an internally calibrated correlation
rather than one adopted from the literature; it is therefore stated in
full in \ref{app:cop} so that the results remain reproducible. Simpler
Lorenz surrogates, and their accuracy relative to detailed cycle models,
have been compared systematically for energy-planning applications
\cite{Pieper.2020}, and the broader consequences of simplifying heat
pump performance models for optimisation results have been examined
directly \cite{Zaubitzer.2026,Kuosa.2024} -- the internally calibrated
correlation used here is one point in that trade-off space, chosen for
consistency with the notebook calibration rather than for minimal model
complexity. The exergetic efficiency is set to $\eta=0.75$ for
both networks, following the calibration notebook underlying this
correlation; \Cref{tab:T1} lists it alongside the other network-specific
parameters.

Two properties of \cref{eq:cop} govern how it enters the optimisation.
First, the exergetic efficiency $\eta$ is a scalar; the temperature
dependence of $\COP(t)$ arises entirely from the source and sink
temperature series. Second, and more consequentially, $\COP(t)$ is
evaluated \emph{outside} the optimisation, once per transformation
stage, from the exogenous supply temperature series produced by
\cref{eq:heatcurve}, and is passed to the solver as a fixed parameter.
The performance figure is thus consistently propagated into every stage
but is not itself a decision variable, and \cref{eq:pel} remains linear.
We state this explicitly because the coupling between supply
temperature and heat pump performance is central to the results in
Section~\ref{sec:results:coupling}, and that coupling is a modelled
physical consequence of the enumerated stage, not an outcome the solver
negotiates.

Heat pump operation is subject to a commitment binary $u(t)$ and a
minimum part load, giving the standard product of a binary with a
bounded continuous capacity that is treated in
Section~\ref{sec:model:objective}. The electrode boiler is modelled as a
continuous investment capacity
$\dot{Q}_{\mathrm{EB}} \in [0, \overline{Q}_{\mathrm{EB}}]$ with
conversion efficiency $\eta_{\mathrm{EB}}$ (\num{0.98} and \num{0.95}
for the two networks), consistent with reported performance figures for
electrode boilers in district heating power-to-heat applications
\cite{DiBacciCapaci.2022,Bottger.2019}, and is available in every
non-baseline scenario.
It is always sited at the same node as the heat pump; its capacity is
consequently a result rather than a scenario dimension, which is why it
does not appear as a column in the scenario matrix. One network
additionally contains an existing \SI{9.9}{\mega\watt} power-to-heat
boiler at the same node. It carries no investment decision:
there is no build binary and no annualised capital cost attached to it,
only variable operating cost. It is therefore complementary to, and
independent of, the investable electrode boiler.

\subsection{Thermal energy storage model}
\label{sec:model:tes}

Storage capacity is selected from a discrete ladder rather than sized
continuously. The ladder is anchored to hours of mean network load,
$\{0, 1, 2, 4, 6, 8, 12, 16, 24, 36, 48, 96, 168\}$ hours, and is
truncated where it would exceed a maximum vessel volume of
\SI{5000}{\cubic\metre}. This ceiling is derived from thin-wall pressure
vessel sizing at the storage design pressure of \SI{10}{\bar}, with a
shell thickness of approximately \SI{60}{\milli\metre} representing what
is practically constructible under the design and conformity-assessment
regime for EU pressure equipment \cite{PED.2014}; it is applied
identically to both
networks. For Stadtbach, this is an exact $\times\frac{1}{10}$ rescaling
of the pre-existing ladder,
$\{0, 7.3, 14.6, 21.9, 29.2, 36.5, 43.8, 51.1, 58.4, 65.7, 73.0,
84.5\}\,\si{\mega\watt\hour}$; for Memmingen, whose previous ladder was
never volume-bound, the same $\{0,1,2,4,6,8,12,16,24,36,48\}$-hour
anchor set is retained and only its two longest-duration steps are
dropped for exceeding the cap ($96\,\mathrm{h}/\SI{134.4}{\mega\watt
\hour}$ and $168\,\mathrm{h}/\SI{235.2}{\mega\watt\hour}$), giving
$\{0, 1.4, 2.8, 5.6, 8.4, 11.2, 16.8, 22.4, 33.6, 50.4,
67.2\}\,\si{\mega\watt\hour}$. Anchoring the ladder to load rather than to
catalogue vessel sizes keeps the two cases comparable despite a
difference in scale of more than a factor of fifty. The chosen minimum
supply--return temperature spread of \SI{15}{\kelvin}, which is enforced
as a hard floor throughout, was confirmed to keep both networks' ladders
within the volume cap even in the worst case
-- Memmingen's top rung reaches \SI{3962}{\cubic\metre}, a margin of
\SI{1038}{\cubic\metre} (\SI{21}{\percent}); Stadtbach's top rung
reaches \SI{4982}{\cubic\metre}, a margin of only \SI{18}{\cubic\metre}
(\SI{0.4}{\percent}), tight by construction rather than coincidence
since it is a rescaling of a ladder originally built to sit exactly at
the previous, larger cap -- so the temperature floor and the volume cap do
not conflict.

Selection is exact:
\begin{equation}
\label{eq:tesladder}
E^{\max} = \sum_{s \in \mathcal{S}} E_s\, z_s,
\qquad \sum_{s \in \mathcal{S}} z_s = 1,
\qquad z_s \in \{0,1\},
\end{equation}
so that capacity is linear in the selection binaries and no big-M
coupling is required. This formulation replaced an earlier one in which
volume was continuous and linked to an activation binary through a big-M
constraint. The replacement is reported here because it is a
transferable modelling result rather than an implementation detail: the
weak linear relaxation of the big-M variant left precisely those
scenarios that place storage at a consumer node stalled at MIP gaps
between \SI{43}{\percent} and \SI{88}{\percent}, and these are the
scenarios that turn out to carry the strongest economic effect
(Section~\ref{sec:results:siting}). A formulation choice that appears
innocuous can therefore suppress the very result the study is designed
to find.

Energy capacity is converted to volume through a stage-specific
coefficient $\xi(\mathrm{HK})$ that depends on the temperature spread
available at that transformation stage, $V = E^{\max} / \xi(\mathrm{HK})$.
An immediate consequence, and one that is easily overlooked -- including
in how results are reported, not only in the formulation itself
(Sections~\ref{sec:results:sweep}--\ref{sec:results:costs}) -- is that
the same rung of the energy ladder corresponds to a different physical
vessel volume at each stage: lowering the supply temperature reduces the
usable spread and therefore inflates the volume, and hence the capital
cost, required to store a given quantity of heat. This is a direct
coupling between the transformation stage and the storage investment,
and it operates in the opposite direction to the efficiency gain that
the same stage confers on the heat pump. The two effects are what the
co-design question in this paper is ultimately about.

Volumes exceeding a standard unit vessel of \SI{25000}{\cubic\metre}
would be realised as $N = \lceil V / V_{\mathrm{unit}} \rceil$ identical
tanks, with the fixed cost incurred per tank rather than per capacity
step, preserving the economies of scale of tank storage while respecting
constructional limits \cite{Dahash.2021}. Under the revised volume cap
of \SI{5000}{\cubic\metre} (\cref{sec:model:tes} above), however, no
scenario in either network can reach the unit-vessel threshold, so the
multi-tank mechanism is present in the formulation but not exercised by
any of the results reported here -- the same status as the multi-pipe
node-mixing constraint of Section~\ref{sec:model:objective}. We report
it for completeness and because it would activate at a higher pressure
rating or a relaxed vessel-sizing assumption. The state of charge
evolves as
\begin{equation}
\label{eq:soc}
E(t) = \lambda^{\Delta t} E(t-1)
 + \eta^{\mathrm{ch}} q^{\mathrm{ch}}(t)\,\Delta t
 - \frac{q^{\mathrm{dis}}(t)\,\Delta t}{\eta^{\mathrm{dis}}},
\qquad 0 \le E(t) \le E^{\max},
\end{equation}
with hourly retention factor $\lambda = \num{0.9999}$, corresponding to a
self-discharge of \SI{0.01}{\percent} per hour, and charging and
discharging efficiencies of \num{0.98} in both networks -- typical
figures for large sensible hot-water storage
\cite{Dahash.2021,DEA.2024b}. The initial
state of charge is fixed at half the installed capacity, and a terminal
state of charge of at least half the installed capacity is enforced as an
inequality. This prevents end-of-horizon depletion without imposing
strict cyclicity: the storage is permitted to finish the year fuller than
it began, but not emptier.

Siting is treated in three variants. In the fixed variants the storage
node is prescribed from a set of candidate nodes; in the endogenous
variant a binary siting vector is optimised over the same candidate set;
and in the colocation variant storage is constrained to share the node
of the heat pump. Comparing the three is the subject of
Section~\ref{sec:results:siting}.

\subsection{Investment layer, objective function and linearisation}
\label{sec:model:objective}

Capital expenditure is annualised following VDI~2067 \cite{VDI2067.2012},
\begin{equation}
\label{eq:annuity}
\mathrm{ANF}(i,n) = \frac{i\,(1+i)^{n}}{(1+i)^{n} - 1},
\end{equation}
with discount rate $i = \num{0.05}$ and component lifetimes of
\num{20}, \num{25} and \num{30} years for the heat pump, the electrode
boiler and the storage respectively \cite{DEA.2024a,DEA.2024b}.
Investment cost is linear with a
fixed activation component,
\begin{equation}
\label{eq:capex}
\mathrm{CAPEX} = \alpha_{\mathrm{HP}} \QWP + \beta_{\mathrm{HP}} y_{\mathrm{HP}}
+ \alpha_{\mathrm{EB}} \dot{Q}_{\mathrm{EB}} + \beta_{\mathrm{EB}} y_{\mathrm{EB}}
+ \alpha_{\mathrm{TES}} V + \beta_{\mathrm{TES}} N,
\end{equation}
and the objective is total annualised cost,
\begin{equation}
\label{eq:tac}
\TAC = \sum_{c} \mathrm{ANF}(i, n_c)\, \mathrm{CAPEX}_c + \mathrm{OPEX},
\end{equation}
where OPEX comprises electricity procurement including the demand
charge, fuel, carbon cost, and a curtailment penalty on
$Q^{\mathrm{dump}}$.

Two distinct linearisation techniques appear in the formulation, and
conflating them would misstate the model's error properties.

The first is a McCormick relaxation \cite{McCormick.1976,Castro.2015} of the bilinear
products of a mass flow with a temperature, $W = \dot{m}\, c_p\, T$,
both factors being continuous variables. For bounds
$\dot{m} \in [\underline{m}, \overline{m}]$ and
$T \in [\underline{T}, \overline{T}]$ the convex hull is
\begin{equation}
\label{eq:mccormick}
\begin{aligned}
W &\ge c_p\left(\underline{m} T + \dot{m}\underline{T} - \underline{m}\,\underline{T}\right), &
W &\ge c_p\left(\overline{m} T + \dot{m}\overline{T} - \overline{m}\,\overline{T}\right), \\
W &\le c_p\left(\overline{m} T + \dot{m}\underline{T} - \overline{m}\,\underline{T}\right), &
W &\le c_p\left(\underline{m} T + \dot{m}\overline{T} - \underline{m}\,\overline{T}\right).
\end{aligned}
\end{equation}
Seven such products occur: pipe inlet and outlet enthalpy, mixing at
nodes with more than one incoming pipe, consumer withdrawal on the
supply and on the return side, and local generation feed-in on the
supply and on the return side. \Cref{eq:mccormick} is an outer
approximation and therefore introduces a relaxation gap, quantified in
Paper~1 \cite{Ruess.2026}; tighter, piecewise variants exist for the
same class of heat-and-electricity bilinear products
\cite{Deng.2021,Castro.2015} but are not required at the problem sizes
solved here. Two of the seven are inactive in the present
study: both networks are tree topologies, so no node receives more than
one incoming pipe and the mixing constraint is never binding, and return
temperature is a stage-fixed parameter rather than a variable in both
configurations. We report them nonetheless, since the constraints are
part of the formulation and would activate on a meshed topology.

The second technique concerns the product of a commitment binary with a
bounded continuous capacity, $\QWP u(t)$. This is reformulated exactly,
\begin{equation}
\label{eq:bigm}
q(t) \le \overline{Q}\, u(t), \quad
q(t) \le \QWP, \quad
q(t) \ge \QWP - \overline{Q}\left(1 - u(t)\right), \quad
q(t) \ge 0,
\end{equation}
which is tight at every integer-feasible point and contributes no
relaxation error. Two further products are not bilinear at all and
require no treatment: $\dot{Q}/\COP$ in \cref{eq:pel}, because $\COP$ is
a parameter, and the energy-to-volume conversion in
Section~\ref{sec:model:tes}, because the temperature spread is fixed
within a stage.

One departure from the $\Lplus$ formulation of Paper~1 was necessary for
tractability. Applying full McCormick temperature propagation across all
\num{8760} hours proved computationally intractable for the smaller
network at the scenario counts required here; a linear supply
temperature offset per pipe is used in its place. The consequences for
the spatial temperature profile are examined in
Section~\ref{sec:discussion:limitations}.

\subsection{Validation strategy}
\label{sec:model:validation}

Five checks are applied and reported in \cref{tab:T5}.

\begin{table}[htbp]
  \centering
  \caption{Model validation summary across the campaign.}
  \label{tab:T5}
  \begin{tabular}{lll}
  \toprule
  Validation metric & Value & Note \\
  \midrule
  Per-node energy-balance agreement & \placeholder{pending} & against live solver state \\
  MIP gap (median, 46 canonical runs) & \placeholder{pending} & enumeration sub-pairs excluded \\
  MIP gap (max, 46 canonical runs) & \placeholder{pending} & \\
  Solver status (46 canonical runs) & \placeholder{pending} & \\
  Annual COP range & \placeholder{pending} & plausibility band \\
  Paper~1 consistency (Memmingen) & \placeholder{pending} & tolerance-band check, Section~\ref{sec:model:validation} \\
  Sweep--MILP optimum consistency & \placeholder{pending} & Section~\ref{sec:results:sweep} \\
  \bottomrule
  \end{tabular}
\end{table}
\placeholder{Do NOT populate this table from any campaign run predating
the eta/CAPEX/TES parameter fixes and the T5 canonical-scenario filter
fix (Section~\ref{sec:results}, status header) -- an earlier campaign
snapshot exists in the project files but is stale on all three counts
and must not be transcribed here.}

Energy conservation is verified by a per-node audit that reads the
solved state directly and compares total heat entering each node against
demand plus onward delivery. Agreement is within
\SI{0.04}{\percent}. We report this rather than an aggregate
generation-plus-net-storage-against-demand statistic, because the latter
is not a valid conservation test in the presence of storage that cycles
within the hour: netting charge against discharge collapses the
throughput of a storage unit that has moved substantial energy in both
directions to approximately zero, and the resulting apparent residual
reflects the arithmetic of the test, not a violation of
\cref{eq:heatbalance}. The distinction matters here because the
best-performing configuration is precisely one in which the storage
cycles in the great majority of hours.

Solver convergence is reported as gap statistics over the
\num{46} canonical scenario runs. Auxiliary runs used to enumerate
siting pairs are excluded from this population and reported separately,
since they are deliberately solved to a fixed time budget rather than to
proven optimality.
\placeholder{Cross-network convergence asymmetry, disclose explicitly
rather than let a reviewer find it: as of this pass, Stadtbach's
scenarios converge to proven optimality (0\% gap) while Memmingen's do
not close within a 24-hour budget (${\sim}$2.5\% gap). Gaps in this
range are standard for MILP energy-system investment models of this
size and are not, by themselves, a problem -- but because the paper's
central claim is a head-to-head comparison between exactly these two
networks, an unexplained asymmetry invites the question of whether part
of the apparent siting difference is numerical slack rather than a real
topology effect. Before writing Section~\ref{sec:results:siting}: (i)
check whether any headline TAC comparison changes conclusion within a
${\sim}$3\% margin -- if none do, state that explicitly; (ii)
additionally check, for the endogenous-siting scenarios specifically
(S6/S7 Stadtbach, S4/S5 Memmingen), whether the \emph{chosen node}
would still be chosen if the remaining gap were resolved in the least
favourable direction -- a stable TAC value does not by itself guarantee
a stable siting decision; (iii) state the asymmetry and both checks
together in this section, using the same disclosure discipline already
applied to the $\eta$/CAPEX harmonisation
(Section~\ref{sec:cases:parameters}).}

Paper~1 validates the network physics -- hydraulics, losses, transport
delay -- at the dispatch level; the present work validates the
investment and siting layer built on top of it. Consistency between the
two is checked through a dedicated run in which heat pump and electrode
boiler capacity are held at Paper~1's values for this network's real,
already-installed assets (\SI{5}{\mega\watt} heat pump,
\SI{5}{\mega\watt} electrode boiler, \SI{500}{\mega\watt\hour} thermal
storage under a cyclic boundary condition), reported as agreement within
a stated tolerance rather than exact reproduction. This run is a
validation instrument only and is not part of the \num{46}-scenario
matrix; its cyclic storage boundary condition differs from the
terminal-inequality condition used elsewhere in this paper
(Section~\ref{sec:model:tes}), appropriate here since capacity is fixed
and given rather than optimised.
\placeholder{Insert the achieved tolerance once the validation run
completes.}

Two plausibility checks complete the set: annual mean coefficients of
performance are required to fall within a physically defensible band for
the temperature lifts involved, and the discrete optimum is required to
coincide, within one grid step and \SI{10}{\percent} in total annualised
cost, with the minimum of the continuous capacity surface obtained from
the parametric sweep of Section~\ref{sec:results:sweep}.

%% ============================================================
\section{Case studies and scenario design}
\label{sec:cases}
%% BUDGET: ~900 words | ELEMENTS: T1, F2, T2
%% STATUS: WRITABLE NOW -- config-derived, campaign-independent
%% ------------------------------------------------------------
The two networks are selected as a deliberate contrast pair rather than
as a convenience sample. \Cref{tab:T1} summarises their key
characteristics; \Cref{fig:F2} maps both topologies with the storage
siting candidates used in Section~\ref{sec:cases:scenarios}.

\begin{figure}[htbp]
  \centering
  \IfFileExists{fig_F2_network_maps.pdf}
    {\includegraphics[width=0.95\linewidth]{fig_F2_network_maps.pdf}}
    {\fbox{\parbox{0.85\linewidth}{\centering\vspace{2cm}
    fig\_F2\_network\_maps.pdf not found --- see note below.\vspace{2cm}}}}
  \caption{District-heating case-study networks and TES candidate sites.}
  \label{fig:F2}
\end{figure}
\placeholder{The source file
\texttt{paper2\_package/02\_figures\_tables/fig\_F2\_network\_maps.svg}
already exists (Fraunhofer colour scheme, both networks). pdflatex
cannot include SVG directly: convert once, e.g.\
\texttt{inkscape fig\_F2\_network\_maps.svg -{}-export-type=pdf} (or
rsvg-convert), then copy the resulting
\texttt{fig\_F2\_network\_maps.pdf} into the same folder as this
\texttt{.tex} file so the \texttt{\textbackslash IfFileExists} check
above picks it up automatically -- no further edit needed here. Also
note the SVG's own title text currently reads ``Stadtbach -- 32 nodes,
meshed''; both ``32'' (Table~T1 says 33) and ``meshed'' (corrected to
``tree'' in Section~\ref{sec:cases:networks}) should be fixed in the
figure-generation script before re-export, not just in this document.}

\subsection{Network characteristics}
\label{sec:cases:networks}

\begin{table}[htbp]
  \centering
  \caption{Key characteristics of the two case-study networks.}
  \label{tab:T1}
  \begin{tabular}{lrr}
  \toprule
  Characteristic & Stadtbach & Memmingen \\
  \midrule
  Nodes (total) & 33 & 15 \\
  \quad producer / mixed / consumer / junction & 6 / 0 / 24 / 3 & 0 / 2 / 13 / 0 \\
  Pipes (supply+return pairs) & 32 & 14 \\
  Total route length [km] & 54.2 & 4.4 \\
  Flow structure & tree, 1 bidirectional trunk & tree, strictly radial \\
  Demand zones & 24 & 27 \\
  Peak heat load [MW] & 200.7 & 4.9 \\
  Annual heat supplied [GWh] & 640.4 & 9.5 \\
  Design supply temp. [\si{\celsius}] & 99.0 & 86.0 \\
  Return temp. [\si{\celsius}] & 60.0 & 63.6 \\
  Fixed generators & 7 & 3 \\
  Investable assets & HP, EB, TES & HP, EB, TES \\
  \bottomrule
  \end{tabular}
\end{table}

Stadtbach is a \num{33}-node tree of \num{32} pipes -- one of them,
the trunk, not constrained to a fixed flow direction -- connecting a
\SI{357.5}{\mega\watt} fixed generation portfolio -- gas- and
biomass-fired combined heat and power, gas boilers and an existing
waste-heat feed -- to \num{24} consumer nodes at a peak load of
\SI{200.7}{\mega\watt} and an annual supply of \SI{640.4}{\giga\watt
\hour}. \placeholder{Confirm pipe count (32) and single connected
component.} Its design supply temperature of \SI{99}{\celsius}
leaves substantial headroom for heat curve reduction, and the
non-fixed flow direction on its trunk pipe makes siting a genuinely
open spatial question rather than one with an obvious answer along a
single branch.

Memmingen is the same network used as the primary case in Paper~1
\cite{Ruess.2026}: a strictly radial \num{15}-node tree of \num{14}
pipes, every pipe fixed to a single flow direction, serving
\num{27} consumers across \num{3} fixed generators, at a peak load of
\SI{4.9}{\mega\watt} against an annual supply of \SI{9.5}{\giga\watt
\hour} -- roughly forty times smaller in peak load and sixty times
smaller in annual demand than Stadtbach; these figures supersede those
given for the same site in Paper~1. Its lower design
supply temperature of \SI{86}{\celsius} and its \num{27} demand zones,
aggregated onto fewer physical nodes than Stadtbach, complete the
picture.

The pairing is deliberate on geometric grounds. Stadtbach's
non-fixed-direction trunk and larger candidate-node set give the siting
problem real spatial freedom to exploit; Memmingen's strictly radial,
single-direction topology and an order of magnitude smaller scale
represent the opposite end of that spectrum, where a network's shape
leaves siting comparatively little room to matter. Both networks are
optimised against the same cost-minimising objective
(Section~\ref{sec:model:objective}) using real, metered operating data;
no asymmetry in economic driver is assumed or intended between them.
Any finding that holds across both networks despite this geometric
contrast is therefore attributable to topology and scale rather than to
a single network's idiosyncrasies. Section~\ref{sec:results:siting}
exploits precisely this contrast: the two networks turn out to favour
different siting strategies, and the geometric characteristics
summarised here are the basis on which that difference is interpreted.

\subsection{Scenario matrix}
\label{sec:cases:scenarios}

\begin{table}[htbp]
  \centering
  \footnotesize
  \caption{Scenario matrix (16 families, 46 MILP runs). Heat-curve
  stages: TVLFIX (constant $T_\mathrm{s}$), HK0/HK1/HK2 ($\kHC=1.0/0.8/0.6$).}
  \label{tab:T2}
  \resizebox{\textwidth}{!}{%
  \begin{tabular}{llllllll}
  \toprule
  Case & Family & Description & TES node & HP/EB node & HK stages & Runs & Features \\
  \midrule
  Stadtbach & BC & Baseline, no investment & --- (no TES) & j\_man & TVLFIX & 1 & --- \\
  Stadtbach & S0 & HP+EB investment, no TES & --- (no TES) & j\_man & TVLFIX, HK0, HK1, HK2 & 4 & --- \\
  Stadtbach & S1 & TES at central node & j\_hkw & j\_man & HK0, HK1, HK2 & 3 & --- \\
  Stadtbach & S2 & TES at HP node & j\_man & j\_man & HK0, HK1, HK2 & 3 & --- \\
  Stadtbach & S3 & TES at HP node, hot charging & j\_man & j\_man & HK0, HK1, HK2 & 3 & hot charging \\
  Stadtbach & S4 & TES at pumping station A & j\_pss & j\_man & HK0, HK1, HK2 & 3 & --- \\
  Stadtbach & S5 & TES at pumping station B & j\_psw & j\_man & HK0, HK1, HK2 & 3 & --- \\
  Stadtbach & S6 & Free HP+TES siting & solver-chosen (5 candidates) & solver-chosen & HK0, HK1, HK2 & 3 & endogenous siting \\
  Stadtbach & S7 & Free siting, co-located, hot charging & solver-chosen (5 candidates) & = TES & HK0, HK1, HK2 & 3 & hot charging, endogenous, co-located \\
  Memmingen & BC & Baseline, no investment & --- (no TES) & j\_12 & TVLFIX & 1 & --- \\
  Memmingen & S0 & HP+EB investment, no TES & --- (no TES) & j\_12 & TVLFIX, HK0, HK1, HK2 & 4 & --- \\
  Memmingen & S1 & TES at central node & j\_1 & j\_12 & HK0, HK1, HK2 & 3 & --- \\
  Memmingen & S2 & TES at HP node & j\_12 & j\_12 & HK0, HK1, HK2 & 3 & --- \\
  Memmingen & S3 & TES at HP node, hot charging & j\_12 & j\_12 & HK0, HK1, HK2 & 3 & hot charging \\
  Memmingen & S4 & Free HP+TES siting & solver-chosen (6 candidates) & solver-chosen & HK0, HK1, HK2 & 3 & endogenous siting \\
  Memmingen & S5 & Free siting, co-located, hot charging & solver-chosen (6 candidates) & = TES & HK0, HK1, HK2 & 3 & hot charging, endogenous, co-located \\
  \bottomrule
  \end{tabular}}
\end{table}
\placeholder{Node codes (j\_hkw, j\_man, j\_pss, j\_psw, j\_1, j\_12) are
internal config identifiers -- decide whether to keep them as-is, map
them to the more descriptive station names used in prose
(Section~\ref{sec:cases:networks}), or move this table to the
supplement given its width; it counts as one of the 12 main-text
elements against the 15-element ECM limit either way
(Section~\ref{sec:cases}, element budget tracker at the end of this
file).}

Each network is subjected to the same scenario design, adapted to its
number of storage siting candidates. A baseline (BC) fixes heat pump,
electrode boiler and storage capacity at zero and a constant supply
temperature, establishing a no-flexibility reference point rather than
necessarily the literal present-day system.
\placeholder{For Memmingen this distinction matters: its
real, currently installed system already includes a 5\,MW heat pump,
5\,MW electrode boiler and 500\,MWh storage (Paper~1,
Section~\ref{sec:model:validation}), so BC for that network is a
counterfactual zero-flexibility case, not a description of what is
actually built today. State this explicitly wherever BC is first
introduced, so a reader does not read the reported BC-to-best-scenario
savings as ``what electrification would add to nothing,'' when for this
network it is actually being compared against removing existing assets
first.} A
no-storage investment case (S0) establishes what heat pump and electrode
boiler investment alone achieve, evaluated at the fixed supply
temperature and at all three transformation stages (Section
\ref{sec:model:heatcurve}). A sequence of fixed-siting scenarios then
places storage at each of several candidate nodes in turn -- the central
generation hub, the node that hosts the heat pump, and, for
Stadtbach, two further pumping-station nodes -- crossed
with the three transformation stages. One fixed-siting variant per
network additionally enables \emph{hot charging}: the heat pump is
permitted to charge the storage above the network supply temperature,
decoupling the usable energy density of the tank from the network's own
temperature spread at the cost of an additional charging-side
temperature lift. Finally, an endogenous-siting pair of scenarios lets
the solver choose the storage node freely from the same candidate set --
once with the heat pump held at its fixed node, and once with
the heat pump siting itself co-optimised and constrained to coincide
with the storage node, again with hot charging enabled. This yields
\num{26} scenario runs for Stadtbach and \num{20} for
Memmingen, \num{46} in total, each solved as an independent
MILP; the two endogenous-siting pairs additionally require the
enumeration procedure described in Section~\ref{sec:model:validation}
where the monolithic formulation does not converge.

The design isolates three questions in a controlled way: what fixed-
siting comparison identifies as the best available node
(Section~\ref{sec:results:costs}), whether free siting improves on that
best fixed choice (Section~\ref{sec:results:siting}), and how each of
these interacts with the transformation stage
(Section~\ref{sec:results:coupling}).

\subsection{Economic and technical parameters}
\label{sec:cases:parameters}

Capital cost is annualised at a discount rate of \SI{5}{\percent}
\cite{DEA.2024a} with
component lifetimes of \num{20}, \num{25} and \num{30} years for heat
pump, electrode boiler and storage respectively \cite{DEA.2024a,DEA.2024b}
(\cref{eq:annuity}).
Unit investment costs $\alpha$ and activation costs $\beta$
(\cref{eq:capex}) are literature- and VDI-2067-style assumptions rather
than vendor quotations, and are compiled in \cref{tab:T1}
\placeholder{add a compact cost-parameter table or fold into T1 -- values
below are ready to insert}. The
heat pump unit cost, $\alpha_{\mathrm{HP}} = \SI{700000}{\eur\per
\mega\watt}$, is applied identically to both networks and is supported
by a survey of 26 realised large-scale heat pump installations in Danish
district heating, which reports specific investment costs of
\SIrange{0.8}{1.1}{\mega\eur\per\mega\watt} \cite{Pieper.2018}, corroborated
by IRENA and JRC figures of the same order. The electrode boiler unit
cost, $\alpha_{\mathrm{EB}} = \SI{150000}{\eur\per\mega\watt}$ with
activation cost $\beta_{\mathrm{EB}} = \SI{20000}{\eur}$, and the storage
unit cost, $\alpha_{\mathrm{TES}} = \SI{1200}{\eur\per\cubic\metre}$
(harmonised, both networks, following the pressure-vessel-sizing
revision in Section~\ref{sec:model:tes}), follow the same Danish Energy
Agency technology data \cite{DEA.2024a,DEA.2024b}.
\placeholder{Insert the current $\beta_{\mathrm{TES}}$ activation cost
for each network here -- not confirmed whether this was harmonised
alongside alpha\_TES or retained at its earlier per-network split.
Confirm the exact current-edition figures for HP/EB
lifetime and TES lifetime/cost against the live DEA datasheets before
submission -- the citation type is correct and standard for this
literature, but the precise numbers (20/25/30 years; the EB and TES
alpha/beta values) were not independently verified against the current
catalogue edition in this pass, only shown to be the right order of
magnitude and the right kind of source.}
Their influence on the
headline results is bounded by the sensitivity analysis in
Section~\ref{sec:discussion:sensitivity}. Operating cost is driven by
network-specific hourly time series for electricity price
\cite{ENTSOE.2025}, grid carbon
intensity \cite{ElectricityMaps.2025}, ambient temperature and, for Stadtbach, waste-
heat availability and temperature, all calibrated following the
procedure established in Paper~1 \cite{Ruess.2026}
(Section~\ref{sec:model:framework}) and, for Memmingen, simplified in
the same way as Stadtbach's (Section~\ref{sec:cases:networks}).
\placeholder{Confirm country, year and interval for the
ElectricityMaps.2025 entry below to match what was actually pulled --
currently assumed Germany, hourly, and the model's reference year.}

%% ============================================================
\section{Results}
\label{sec:results}
%% BUDGET: ~2,300 words | ELEMENTS: F3, F4, F5, F6, T3, T4
%% STATUS: BLOCKED -- as of 2026-07-20, the ENTIRE 46-scenario matrix for
%%   BOTH networks is being re-solved following THREE harmonised parameter
%%   fixes, applied in sequence: (1) HP efficiency eta=0.75 for both
%%   networks, (2) HP unit CAPEX 700k EUR/MW for both networks
%%   (literature-supported, see Section 3.3), (3) TES volume cap tightened
%%   50,000->5,000 m3 (ASME/PED thin-wall pressure vessel derivation) and
%%   unit cost raised 400->1,200 EUR/m3 (3x) for both networks -- flagged
%%   as the single largest fix of the day, expected to shift optimal TES
%%   sizing across most scenarios in both networks, not merely equalise
%%   them. Do NOT draft against any KPI number seen in chat before this
%%   date -- including the SB-S2-HK2 headline figure quoted at the start
%%   of the drafting session. A process-management bug (zombie chains from
%%   incomplete kills) was also found and fixed; all four campaigns were
%%   confirmed relaunched cleanly (5th generation) after a full sweep.
%%   F3 sweep and F7 tornado need fresh launches for BOTH networks once
%%   the four waves land -- the ladder and cost changes above invalidate
%%   any sweep grid computed under the old TES parameters too. T5
%%   exclusion filter logic confirmed already correct in code; still needs
%%   re-verification against fresh output. SB-S2 plausibility check
%%   (Part G.4, TES-at-consumer-node economics) still outstanding and now
%%   applies to whatever scenario turns out to be the new Stadtbach
%%   headline under the revised TES economics.
%%   CONVERGENCE ASYMMETRY (as of this pass): Stadtbach proven-optimal
%%   (0%), Memmingen ~2.5% after 24h, not closing further. Acceptable by
%%   domain norms, but the two required checks in Section 2.7's
%%   validation placeholder (headline-TAC margin, siting-decision
%%   stability under worst-case gap resolution) MUST be done and their
%%   result stated in this section before the siting comparison
%%   (Section 4.4) is written -- do not silently drop this once campaigns
%%   finish, it is not resolved just because the numbers arrive.
%% ------------------------------------------------------------
\subsection{Capacity--cost surface}
\label{sec:results:sweep}
\todo{F3 -- 7x7 dispatch-only LP sweep per network, LCOH colour scale,
MILP optimum marked. Also serves as cross-validation.
REQUIREMENT: label the capacity axis in both MWh and the corresponding
m$^3$ at that scenario's $\Delta T$ (Section~\ref{sec:model:tes}), or
show volume as a secondary axis -- the sweep grid is built around a
fixed MILP optimum whose volume differs by HK stage even at equal MWh,
and an MWh-only axis would obscure that.}

\subsection{Coupling of heat curve, COP and storage capacity}
\label{sec:results:coupling}
\todo{F4 -- THE central figure of the paper. Ceteris-paribus subset:
same siting, HK stage varied only.
CRITICAL REQUIREMENT, do not skip: $V_\mathrm{TES}$ is selected from an
energy ladder (MWh) that is identical across HK stages, but the same
MWh rung maps to a different physical volume -- and hence a different
installed cost -- at each stage, because the energy-to-volume
coefficient $\xi(\mathrm{HK})$ scales with $\Delta T$
(Section~\ref{sec:model:tes}). If this figure plots storage capacity in
MWh alone, that cost penalty is real in the optimisation but invisible
on the page, and the coupling this figure exists to show would silently
disappear. Plot volume (m$^3$) as the storage-axis quantity, or MWh with
an explicit secondary axis/annotation showing the implied volume or
\si{\eur}/MWh at each stage -- confirmed with the author as a hard
requirement, not a style preference.}

\subsection{Cost and emission reduction versus baseline}
\label{sec:results:costs}
\todo{F5 + T3/T4. Stadtbach best fixed siting SB-S2-HK2; Memmingen best
overall MM-S5-HK2. Consider moving the full 46-row matrices to the
supplement -- see element budget note at the end of this file.
REQUIREMENT: give T3/T4 a $V_\mathrm{TES}$ [m$^3$] column alongside the
existing $E_\mathrm{TES}$ [MWh] column (or an equivalent
\si{\eur}/MWh$_\mathrm{installed}$ figure) -- same reasoning as F4 above.
Without it, a reader comparing the "TES capacity" number across HK0 and
HK2 rows for the same scenario family sees identical MWh figures and
has no way to see that HK2's rung costs more, even though the model
priced that difference in correctly.}

\subsection{Fixed versus endogenous siting}
\label{sec:results:siting}
\todo{F6 -- the asymmetry: free siting does not beat a well-chosen fixed
node in Stadtbach, while colocation beats free siting in Memmingen.}

%% ============================================================
\section{Discussion}
\label{sec:discussion}
%% BUDGET: ~900 words | ELEMENTS: F7, T5
%% ------------------------------------------------------------
\subsection{Sensitivity of the design to economic assumptions}
\label{sec:discussion:sensitivity}
\todo{F7 tornado}

\subsection{Why network topology decides the siting strategy}
\label{sec:discussion:topology}
\todo{The transferable insight. Requires the SB-S2 economic mechanism
from the G.4 plausibility check -- do not write before it lands.}

\subsection{Model validation and limitations}
\label{sec:discussion:limitations}
\todo{T5. Limitations: staged rather than continuous heat curve;
linear temperature offset instead of full McCormick propagation;
cycling not penalised in the objective (state explicitly, and report
the counter-test if one is run); single weather/price year.}

%% ============================================================
\section{Conclusions}
\label{sec:conclusions}
%% BUDGET: ~400 words | ELEMENTS: none
%% ------------------------------------------------------------
\todo{Write after Sec. 4-5. No new results, no hedging.}

%% ============================================================
\section*{Nomenclature}
\todo{Symbols, subscripts, abbreviations -- keep consistent with Paper 1.}

%% ============================================================
\section*{CRediT authorship contribution statement}
\todo{Mirror the Paper 1 statement, adjusted for Paper 2 contributions.}

%% ============================================================
\section*{Declaration of competing interest}
\todo{Standard statement.}

%% ============================================================
\section*{Data availability}
%% ECM applies research-data Option C: deposit + cite, or justify.
\todo{Deposit CALION configs, scenario definitions and KPI outputs in a
repository and cite the DOI here. Plan this BEFORE submission -- it is
a hard requirement for this journal, not a courtesy.}

%% ============================================================
\section*{Acknowledgements}
\todo{Funding sources in the Elsevier standard phrasing.}

%% ============================================================
%% MUST appear directly before the reference list (ECM requirement)
\section*{Declaration of generative AI and AI-assisted technologies in
the manuscript preparation process}
\todo{Declare tool use, incl. the coding agent where it contributed to
figure generation or analysis code. Note: AI-generated data
visualisations are permitted when derived from underlying data by
reproducible methods, but must be disclosed in the figure captions as
well as here.}

%% ============================================================
\appendix

\section{Coefficient of performance correlation}
\label{app:cop}

The coefficient of performance in \cref{eq:cop} is evaluated from the
supply temperature series of the respective transformation stage. With
the logarithmic mean temperature
\begin{equation}
L(T_{\mathrm{h}}, T_{\mathrm{c}})
= \frac{T_{\mathrm{h}} - T_{\mathrm{c}}}{\ln (T_{\mathrm{h}}/T_{\mathrm{c}})},
\end{equation}
the source outlet temperature
$T_{\mathrm{src,out}}(t) = \max\{T_{\mathrm{src}}(t) - \Delta T,\, 1\}$
and the sink outlet temperature
$T_{\mathrm{sink,out}}(t) = T_{\mathrm{s}}(t)$ from
\cref{eq:heatcurve}, define
\begin{align}
L_{\mathrm{sink}}(t) &= L\!\left(T_{\mathrm{sink,out}}(t),\, T_{\mathrm{sink,in}}\right), &
L_{\mathrm{src}}(t) &= L\!\left(T_{\mathrm{src}}(t),\, T_{\mathrm{src,out}}(t)\right).
\end{align}
Two empirical loss terms account for internal temperature differences and
auxiliary consumption,
\begin{align}
m_{\mathrm{dts}}(t) &= 0.2\left(T_{\mathrm{sink,out}}(t) - T_{\mathrm{src,out}}(t) + 2\Delta T_{\mathrm{pp}}\right)
 + 0.2\left(T_{\mathrm{sink,out}}(t) - T_{\mathrm{sink,in}}\right) + 0.016, \\
q_{\mathrm{ww}}(t) &= 0.0014\left(T_{\mathrm{sink,out}}(t) - T_{\mathrm{src,out}}(t) + 2\Delta T_{\mathrm{pp}}\right)
 - 0.0015\left(T_{\mathrm{sink,out}}(t) - T_{\mathrm{sink,in}}\right) + 0.039,
\end{align}
from which the two factors of \cref{eq:cop} follow as
\begin{equation}
A(t) = \frac{L_{\mathrm{sink}}(t)}{L_{\mathrm{sink}}(t) - L_{\mathrm{src}}(t)},
\qquad
B(t) = \frac{1 + \dfrac{m_{\mathrm{dts}}(t) + \Delta T_{\mathrm{pp}}}{L_{\mathrm{sink}}(t)}}
{1 + \dfrac{m_{\mathrm{dts}}(t) + 0.5\left(T_{\mathrm{src}}(t) - T_{\mathrm{src,out}}(t)\right) + 2\Delta T_{\mathrm{pp}}}
{L_{\mathrm{sink}}(t) - L_{\mathrm{src}}(t)}} .
\end{equation}
The result is clipped to the interval $[1.01, 12]$. Parameter values are
$F_Q = \num{0.10}$, $\Delta T = \SI{20}{\kelvin}$ and
$\Delta T_{\mathrm{pp}} = \SI{5}{\kelvin}$ for both networks;
$T_{\mathrm{sink,in}}$ and $\eta$ are network specific and are listed in
\cref{tab:T1}.
\placeholder{Confirm that the logarithmic means are evaluated on absolute
temperatures, and state the calibration basis for the empirical
coefficients -- reviewers will ask where 0.2, 0.016, 0.0014 and 0.039
come from.}

%% ============================================================
\bibliographystyle{elsarticle-num}
\bibliography{Paper2_V1}

\end{document}

%% ============================================================
%% ELEMENT BUDGET TRACKER — ECM allows 15 figures+tables COMBINED
%% ------------------------------------------------------------
%%  Main text (current plan):
%%    F1 model architecture          F2 network maps
%%    F3 capacity heatmap            F4 k/COP/V_TES coupling
%%    F5 cost reduction              F6 siting comparison
%%    F7 sensitivity tornado         T1 network characteristics
%%    T2 scenario matrix             T3 Stadtbach KPIs (reduced)
%%    T4 Memmingen KPIs (reduced)    T5 validation
%%    = 12 elements  -> 3 slots of headroom
%%  Supplement:
%%    F8 spatial T/p profiles, F9 TES SOC time series,
%%    full 46-row KPI matrices, T3b/T4b endogenous siting detail,
%%    T1b generator portfolio
%%  NOTE: the original figure spec (F1-F9 + T1a/T1b/T2-T5) totals 15
%%  BEFORE T3b/T4b, i.e. over budget. The split above resolves it.
%% ============================================================